\scanpage{124}
\mspage{61}
Considérons les images inverses de $T_d$ et $T_{d'}$ sur $D_{d,d'}$, qu'on va noter
\deleted{$\unsure{T_d}_{D_{d,d'}}$}
\[
(193)\quad\left\{\begin{aligned}
T(d,d')_d&=(T_d)_{D_{d,d'}}&&\text{syst. local de rang }d\text{ sur }D_{d,d'}\\
T(\unsure{d|d^*},d')_{d'}&=(T_{d'})_{D_{\unsure{d|d^*},d'}}&&\text{syst. local de rang }d'\text{ sur }D_{d,d'}
\end{aligned}\right.
\]
on a alors un hom. canonique surjectif (dit ``tautologique'')
\[
(194)\qquad T(d,d')_d\longrightarrow T(d,d')_{d'}
\]
qui, par image inverse sur $D'_{d,d'}=D'_d\mathbin\times_{D_d}D_{d'}$ (cf diagramme can.

(195)\quad
\diagram{\begin{tikzcd}[ampersand replacement=\&, column sep=large,row sep=large]
\&\mathcal P_{d'}(I_d)_{D'_d}\simeq D'_{d,d'}\arrow[dl]\arrow[dr,"\unsure{\mathfrak S_d|\mathfrak S_{d'}}"]\arrow[rr,dashed]\&\&D'_{d'}\arrow[dd,"\mathfrak S_{d'}"]\\
D'_d\arrow[dr,dash,"\mathfrak S_d"']\&\&D_{d,d'}\arrow[dl]\arrow[dr]\&\\
\&D_d\&\&D_{d'}
\end{tikzcd}}
s'identifie à l'hom. de projection
\[
(196)\qquad\mathbb Z^d\unsure{\text{---}|\longrightarrow}\mathbb Z^{d'}
\]
de $\mathbb Z^d$ sur un produit partiel défini par une partie de cardinal $d'$ de l'ens $I_d$ en exposant : $\mathbb Z^d\simeq\mathbb Z^{I_d}$.
\margin{123}
